% ==============================================================================
% Resumo na língua vernácula (obrigatório, NBR 14724:2024 e NBR 6028:2021)
% ==============================================================================
\begin{resumo}
\begin{instrucao}
\orientacao{Resumo}
O resumo apresenta, de forma concisa, o objetivo, o método, os resultados e as
conclusões do trabalho. Escreva-o por último, quando o texto estiver pronto.
Siga a NBR 6028:2021:
\begin{itemize}
  \item parágrafo único, com 150 a 500 palavras;
  \item frases afirmativas, com o verbo na terceira pessoa, como recomenda o
        item 4.1.4 da norma (``Este trabalho propõe\ldots'', ``Os resultados
        mostram\ldots''). A voz ativa deixou de ser exigida, mas costuma
        deixar o resumo mais direto;
  \item sem citações, fórmulas, siglas não explicadas, figuras ou tabelas;
  \item palavras-chave logo abaixo, separadas por ponto e vírgula e finalizadas
        por ponto. Prefira termos que um leitor usaria para encontrar o
        trabalho num catálogo (de três a cinco é a prática mais comum).
\end{itemize}
Um roteiro que funciona bem dedica uma ou duas frases ao contexto e ao
problema, outras ao objetivo e ao método e as últimas aos principais
resultados (com números) e à conclusão.
\end{instrucao}

A monitoria acadêmica é uma das principais formas de apoio aos estudantes nas
disciplinas iniciais dos cursos de Computação, mas a marcação de atendimentos
costuma depender de mensagens informais e planilhas, o que gera conflitos de
horário e dificulta o acompanhamento pela coordenação. Este trabalho apresenta
o Monitora, um sistema web para agendamento de monitorias que permite aos
monitores publicar horários de atendimento, aos estudantes reservar esses
horários e aos professores acompanhar a frequência e os temas atendidos. O
sistema foi desenvolvido com um processo iterativo baseado no Scrum, em seis
ciclos de duas semanas. Os requisitos foram levantados em entrevistas com
monitores e professores e especificados por meio de histórias de usuário. A
implementação usa o arcabouço Django, o banco de dados PostgreSQL e
contêineres Docker. A verificação combinou testes automatizados, com cobertura
de 87\% do código, e uma avaliação com 18 usuários, que atribuíram ao sistema
pontuação média de 84,6 na escala de usabilidade de sistemas, próxima da faixa
considerada excelente. Os resultados indicam que o sistema atende aos requisitos
levantados e que pode substituir o processo informal usado atualmente nas
disciplinas do curso.

\palavraschave{sistema web; agendamento; monitoria acadêmica; engenharia de
software}
\end{resumo}
